% Generated by scripts/assemble.py; corrections belong in source/corrections.json.
\examhead{58}{天津大学 2024 年数学分析试卷}{57}
% p058-b18; source PDF page 58
\blockanchor{p058-b18}
\question{1}构造一个二元函数满足以下条件并证明.\par
% p058-b19; source PDF page 58
\blockanchor{p058-b19}
(1) 在原点处连续且偏导数存在;\par
% p058-b20; source PDF page 58
\blockanchor{p058-b20}
(2) 在原点处不可微.\par
% p058-b21; source PDF page 58
\blockanchor{p058-b21}
\question{2}设 $f(x)$ 在 $[a, b]$ 上二阶可导, $f(a) = f(b)$, 且 $|f''(x)| \leq c$, 证明: $|f'(x)| \leq \frac{1}{2} c (b - a)$.\par
% p059-b00; source PDF page 59
\blockanchor{p059-b00}
\question{3}设 $f(x)$ 在 $[a, +\infty)$ 上可导, $f(a) > a$, 且 $f'(x) + f(x) > 1 + x$, 证明: $f(x) > x$.\par
% p059-b01; source PDF page 59
\blockanchor{p059-b01}
\question{4}设 $f(x)$ 在 $[1, +\infty)$ 上一致连续, 证明: 存在 $M > 0$, 满足 $\frac{|f(x)|}{x} \leq M, x \in [1, +\infty)$.\par
% p059-b02; source PDF page 59
\blockanchor{p059-b02}
\question{5}设数列 $\{a_{n}\}$ 满足 $\lim_{n\to \infty}\frac{a_n}{n^\alpha} = 1$ ($\alpha >0$), 求极限 $\lim_{n\to \infty}\frac{a_1 + a_2 + \cdots + a_n}{n^{1 + \alpha}}$.\par
% p059-b03; source PDF page 59
\blockanchor{p059-b03}
\question{6}设 $f(x)$ 是以 $2\pi$ 为周期且在 $[- \pi, \pi]$ 上可积的函数, $a_n, b_n$ 为 $f(x)$ 的傅里叶级数.\par
% p059-b04; source PDF page 59
\blockanchor{p059-b04}
(1) 求 $f(x + 1)$ 的傅里叶系数.\par
% p059-b05; source PDF page 59
\blockanchor{p059-b05}
(2) 求  $f(2x)$  的傅里叶系数.\par
% p059-b06; source PDF page 59
\par\Needspace{6\baselineskip}
\blockanchor{p059-b06}
\question{7}设 $f(x)$ 为 $[a, b]$ 上的单调函数, $g(x)$ 是周期为 $T$ 的连续函数, 且 $\int_0^T g(x) \mathrm{d}x = 0$, 证明:\par
\nopagebreak[4]
% p059-b07; source PDF page 59
\blockanchor{p059-b07}
\displayformula{\lim _ {\lambda \rightarrow + \infty} \int_ {a} ^ {b} f (x) g (\lambda x) \mathrm{d} x = 0.}
% p059-b08; source PDF page 59
\blockanchor{p059-b08}
\question{8}已知序列 $\{x_{n}\}, \{y_{n}\}$ 满足 $x_{n+1} = y_{n} + \frac{1}{2} x_{n}$, 证明: $\{y_{n}\}$ 收敛的充要条件是 $\{x_{n}\}$ 收敛.\par
% p059-b09; source PDF page 59
\par\Needspace{6\baselineskip}
\blockanchor{p059-b09}
\question{9}记 $D_{R} = \{|x| \leq R, |y| \leq R, |z| \leq R \mid R > 0\}$. 求极限\par
\nopagebreak[4]
% p059-b10; source PDF page 59
\blockanchor{p059-b10}
\displayformula{\lim _ {R \rightarrow + \infty} \iiint_ {D _ {R}} (x + y + z) ^ {2} e ^ {- (x ^ {2} + y ^ {2} + z ^ {2})} \mathrm{d} x \mathrm{d} y \mathrm{d} z.}
% p059-b11; source PDF page 59
\blockanchor{p059-b11}
\question{10}求积分 $\int_0^{+\infty} \frac{(\sin x)^2}{x^2} \, \mathrm{d}x$.\par
% p059-b12; source PDF page 59
\par\Needspace{6\baselineskip}
\blockanchor{p059-b12}
\question{11}设 $L$ 为 $x^{2} + y^{2} + z^{2} = 2Rx$ 与 $x^{2} + y^{2} = 2rx$ 的交线 $(R > r > 0)$ 的上半部分 $(z \geq 0)$, 求曲线积分\par
\nopagebreak[4]
% p059-b13; source PDF page 59
\blockanchor{p059-b13}
\displayformula{\int_ {L} \left(y ^ {2} + z ^ {2}\right) \mathrm{d} x + \left(z ^ {2} + x ^ {2}\right) \mathrm{d} y + \left(x ^ {2} + y ^ {2}\right) \mathrm{d} z.}
% p059-b14; source PDF page 59
\blockanchor{p059-b14}
从 z 轴正向看, L 取逆时针方向.\par
% p059-b15; source PDF page 59
\par\Needspace{6\baselineskip}
\blockanchor{p059-b15}
\question{12}设 $\{f_n(x)\}$ 为定义在 $[a, b] \setminus \{x_0\}$ 上的函数列，且在 $[a, b] \setminus \{x_0\}$ 上一致收敛于 $f(x)$，同时\par
\nopagebreak[4]
% p059-b16; source PDF page 59
\blockanchor{p059-b16}
\displayformula{\lim _ {x \rightarrow x _ {0}} f _ {n} (x) = c _ {n} (n = 1, 2, \dots).}
% p059-b17; source PDF page 59
\blockanchor{p059-b17}
(1) 证明: $\lim_{n\to \infty}c_n$ 存在.\par
% p059-b18; source PDF page 59
\blockanchor{p059-b18}
(2) 证明: $\lim_{x \to x_0} f(x)$ 存在, 且 $\lim_{x \to x_0} f(x) = \lim_{n \to \infty} c_n$.\par
